% TOP_PROBLEM: {"id": 3295, "title": "Simultaneous partition of hypergraphs", "category_id": 3, "category": {"id": 3, "name": "graph_theory", "display_name": "Graph Theory", "description": "Problems involving graphs, networks, and their properties.", "slug": "graph-theory", "order_index": 3, "created_at": "2026-07-31T15:26:25.671Z"}, "status": "open", "problem_number": "OPG-46817", "set_id": 12, "statusQualification": "The source little-o notation is made explicit with fixed uniformity and both edge counts tending to infinity."}
% FIRST_POSTED: 2026-10-01
% ENTRY_KIND: statement_only
% UnsolvedMath ID 3295; source catalogue code OPG-46817
% !TeX program = lualatex
% Definition and sources only; this file is not an LLM solution attempt.
\documentclass[11pt,a4paper]{article}
\usepackage[margin=25mm]{geometry}
\usepackage{fontspec}
\setmainfont{lmroman10-regular.otf}[BoldFont=lmroman10-bold.otf,
 ItalicFont=lmroman10-italic.otf,BoldItalicFont=lmroman10-bolditalic.otf]
\usepackage{amsmath,amssymb,mathtools}
\usepackage{unicode-math}
\setmathfont{latinmodern-math.otf}
\AtBeginDocument{\let\setminus\smallsetminus}
\usepackage{xurl,hyperref}
\hypersetup{colorlinks=true,urlcolor=blue}
\setcounter{secnumdepth}{0}
\setlength{\parindent}{0pt}
\setlength{\parskip}{5pt}
\setlength{\emergencystretch}{3em}
\begin{document}
\section*{Simultaneous partition of hypergraphs}\label{3295}
{\small UnsolvedMath ID 3295 · OPG-46817 · Graph Theory · OpenGarden}
\subsection{Definitions and mathematical statement}
Fix an integer $r\ge2$. Let $H_1=(V,E_1)$ and $H_2=(V,E_2)$ be finite
simple $r$-uniform hypergraphs on the same vertex set; each edge is an
$r$-element subset and $m_i=|E_i|$. The edge sets may overlap. For a
partition $\mathcal P=(V_1,\ldots,V_r)$ of $V$ into $r$ labelled,
possibly empty parts, call an edge transversal if it meets every
part, necessarily in exactly one vertex. Let $t_i(\mathcal P)$ count
transversal edges of $H_i$.

The question asks for the following uniform asymptotic assertion:
for every fixed $r$ and every $\varepsilon>0$, is there an integer
$M=M(r,\varepsilon)$ such that whenever $m_1,m_2\ge M$, one partition
$\mathcal P$ simultaneously satisfies
\[
             t_i(\mathcal P)\ge
             \left(\frac{r!}{r^r}-\varepsilon\right)m_i
                         \qquad(i=1,2)?
\]
This spells out the source's $r!m_i/r^r-o(m_i)$ notation in the regime
$\min(m_1,m_2)\to\infty$, with $r$ fixed. No balance condition on the
part sizes or hypothesis on vertex degrees and codegrees is imposed.

Independently placing every vertex in one of $r$ parts with equal
probability makes a given edge transversal with probability $r!/r^r$:
its $r$ vertices must receive all distinct labels. Thus each hypergraph
separately has a partition attaining its own expected count. The
problem is whether one and the same partition can approach both
expectations, even when their edges have very different overlap
patterns and their edge counts grow at different rates.
\subsection{Short English statement}
Can one vertex partition simultaneously retain nearly the random-partition proportion of transversal edges in each of two uniform hypergraphs?
\subsection{Sources}
\begin{itemize}
\item[S1] ULAM.AI and the UnsolvedMath Contributors, supplied problems.json, record 3295 (OPG-46817), OpenGarden collection. Catalogue identity and source transcription; CC BY 4.0.
\url{https://huggingface.co/datasets/ulamai/UnsolvedMath}.
\item[S2] Open Problem Garden, Simultaneous partition of hypergraphs. Original problem and discussion; the mathematical formulation below is expanded from the supplied source record.
\url{http://www.openproblemgarden.org/op/simultaneous_partition_of_hypergraphs}.
\item[S3] D. Kühn and D. Osthus, Maximizing several cuts simultaneously (2005/2007), source paper for the simultaneous-partition question.
\url{https://arxiv.org/abs/math/0503403}.
\end{itemize}
\end{document}
